\section{Rompecabezas conductual: cómo cuantificar el In-Context Learning}

Antes de estudiar los circuitos internos, primero hay que definir el fenómeno externo que se quiere explicar. \term{In-context learning} (ICL) se refiere a que, con los parámetros fijos, el modelo mejora sus predicciones posteriores únicamente a partir de los examples, pattern o instructions presentes en el context actual. Parece como si el modelo ejecutara otra vez un aprendizaje temporal dentro del prompt; por eso es una de las capacidades más llamativas de los modelos de lenguaje grandes.

\lecturefigure{slide-04-in-context-learning-motivation.jpg}{In-context learning: un modelo de lenguaje con parámetros fijos adquiere temporalmente nuevas capacidades a partir del prompt/context.}{Video oficial de Stanford Online 00:05:40--00:06:51.}

\teachervoice{El expositor llama al ICL una de las propiedades más striking de los modelos de lenguaje grandes, pero aquí el «learning» no es una gradient update. Los parámetros del modelo permanecen sin cambios; el cambio ocurre en las residual activations y en la attention computation.}

\subsection{Dos learning axis}

Una learning curve ordinaria toma el avance del entrenamiento como eje horizontal y observa cómo mejora el modelo a medida que aumentan los gradient steps; la ICL curve, en cambio, fija un training snapshot, toma el context token index como eje horizontal y observa si, dentro del mismo forward, los tokens más tardíos se predicen mejor. Al reunir ambas se obtiene una función bidimensional
\[
\mathcal{L}(s,i)
=\mathbb{E}_{x}\left[-\log p_{\theta_s}(x_i\mid x_{<i})\right],
\]
donde:
\begin{itemize}
  \item $s$: training snapshot o tokens de entrenamiento acumulados;
  \item $i$: token index dentro de la sequence actual;
  \item $\theta_s$: parámetros del snapshot $s$;
  \item $\mathcal{L}(s,i)$: next-token loss promedio de ese snapshot en la $i$-ésima context position.
\end{itemize}

\lecturefigure{slide-05-learning-vs-in-context-curves.jpg}{Learning curve e in-context learning curve: eje del tiempo de entrenamiento y eje de la posición en el context.}{Video oficial de Stanford Online 00:06:52--00:10:33.}

\begin{importantbox}{Los dos tipos de «aprendizaje» no deben confundirse}
El training learning modifica $\theta_s$; el in-context learning mantiene fijo $\theta_s$ y solo cambia el prefix de entrada y las activations internas. El primero responde «qué pasa si el modelo se entrena más tiempo»; el segundo responde «qué pasa si el mismo modelo ve más context».
\end{importantbox}

\subsection{Mapa bidimensional de loss e ICL score}

Al apilar la per-token loss de cada snapshot, con la context position en el eje vertical y el snapshot en el eje horizontal, se obtiene un mapa bidimensional. Un corte horizontal es la training curve de una token position; un corte vertical es la ICL curve de un snapshot. El expositor usa un summary statistic sencillo:
\[
S_{\mathrm{ICL}}(s)
=\mathcal{L}(s,500)-\mathcal{L}(s,50).
\]

\begin{itemize}
  \item Si $S_{\mathrm{ICL}}<0$, el token 500 es más fácil de predecir que el token 50;
  \item cuanto más negativo es el valor, más fuerte es la late-context advantage;
  \item este indicador solo condensa la context-position improvement y no equivale a la few-shot task accuracy;
  \item los tokens 50 y 500 son una operational choice de la configuración del artículo, no la única definición legítima de ICL.
\end{itemize}

\lecturefigure{slide-06-token-snapshot-loss-map.jpg}{Gráfica de diagnóstico unificada: loss por token index × training snapshot, y $\mathrm{Loss@500}-\mathrm{Loss@50}$.}{Video oficial de Stanford Online 00:10:34--00:12:45.}

\begin{knowledgebox}{Lectura de la figura: primero los ejes, luego la curva de diferencias a la derecha}
El mapa de calor de la izquierda muestra la loss de cada training snapshot en distintos token index; Loss@50 y Loss@500, al centro, son dos cortes horizontales; la curva de diferencias de la derecha elimina la mejora global del entrenamiento y resalta «si la segunda mitad del context se vuelve útil de repente».
\end{knowledgebox}

\begin{lstlisting}[caption={Cálculo del ICL score y del phase-change diagnostic a partir de la per-token loss.}]
def in_context_learning_score(loss_by_snapshot_and_token):
    loss_50 = loss_by_snapshot_and_token[:, 49]
    loss_500 = loss_by_snapshot_and_token[:, 499]
    score = loss_500 - loss_50  # more negative means stronger ICL
    slope = derivative(score, with_respect_to="log_training_tokens")
    return score, slope
\end{lstlisting}

\subsection{Qué tan grande es realmente 0.4 nats}

Cuando la negative log-likelihood usa el natural logarithm, su unidad es el nat. Para convertirla a bits, se divide entre $\ln 2$:
\[
0.4\ \mathrm{nats}
=\frac{0.4}{\ln 2}\ \mathrm{bits}
\approx 0.577\ \mathrm{bits/token}.
\]
Dicho de otro modo, la uncertainty del modelo sobre los late-context tokens disminuye aproximadamente medio bit por token, lo que equivale aproximadamente a una binary choice menos cada dos tokens. Esto no significa «la exactitud mejora 40\%» ni puede compararse directamente entre conjuntos de datos, pero basta para mostrar que la información nueva dentro del context se aprovecha de forma estable.

\lecturefigure{slide-07-meaning-of-point-four-nats.jpg}{Intuición del tamaño del efecto: $0.4$ nats equivalen a cerca de $0.5$ bits/token.}{Video oficial de Stanford Online 00:12:46--00:13:47.}

\teachervoice{Olah convierte por iniciativa propia los nats en bits para evitar que se subestimen los decimales de la curva. En language modeling cada token acumula pérdida, y medio bit/token en secuencias largas produce una total log-likelihood difference muy grande.}

\subsection{Resumen de la sección}

La métrica conductual del ICL proviene de una loss surface bidimensional: un eje es el avance del entrenamiento y el otro es la context position. $\mathrm{Loss@500}-\mathrm{Loss@50}$ condensa la late-token advantage en una sola curva y proporciona una magnitud objetivo, susceptible de intervención, para la búsqueda posterior del circuit interno.
